\chapter{Data: Schemas, Deal Types, and Point-in-Time Engineering}
\label{ch:data}

Every modelling decision in Parts~I--III is downstream of two facts about the data: the
PitchBook asset is a \emph{transaction log with attached snapshots}, not a panel; and the
FactSet feed is \emph{published with lags}, so the value that existed on a date is not the
value that was knowable on that date. This chapter fixes the schemas, formalizes the
point-in-time discipline that governs every join in the manual, and quantifies the defects
(missingness, noise, duplication) that each later chapter must survive. Readers who skip
it will misread at least the leakage warnings of Chapters~\ref{ch:success}
and~\ref{ch:xgb} and the deduplication conventions of Chapter~\ref{ch:hawkes}.

\section{The PitchBook asset: three tables}
\label{sec:data-pitchbook}

The vendor delivers three relations. The \texttt{companies} table has one row per firm in
the \emph{tracked universe}; its columns are \emph{current} snapshots, refreshed by the
vendor, with no history. The \texttt{deals} table has one row per transaction and is the
only genuinely dated object: everything we treat as a time series of firm state is derived
from it. The \texttt{investors} relation links firms to investors with a start date. The
join keys are \texttt{company\_id} (all three tables), \texttt{deal\_id}, and
\texttt{investor\_id}. Table~\ref{tab:data-fields} groups the fields by function rather
than by table, because the functional grouping is what determines how each field may be
used: identifiers key the joins, deal terms and capital-structure fields are legitimate
dated observations, financials-at-deal are noisy state measurements, and firmography is a
mix of static facts and forbidden end-of-sample snapshots
(Section~\ref{sec:data-pit}).

\begin{table}[htbp]
\centering
\small
\begin{tabular}{lp{9.6cm}}
\toprule
Group (table) & Fields \\
\midrule
Identifiers &
\texttt{company\_id}, \texttt{deal\_id}, \texttt{investor\_id},
\texttt{company\_name}, \texttt{investorname}, \texttt{deal\_date},
\texttt{investorsince} \\
Industry, universe (deals) &
\texttt{primary\_industry\_sector}, \texttt{primary\_industry\_group},
\texttt{primary\_industry\_code}, \texttt{company\_universe},
\texttt{current\_financing\_status} \\
Deal terms (deals) &
\texttt{deal\_type}, \texttt{deal\_type2}, \texttt{deal\_type3}, \texttt{deal\_class},
\texttt{deal\_status}, \texttt{deal\_size}, \texttt{deal\_size\_status},
\texttt{pre\_money\_valuation}, \texttt{post\_valuation},
\texttt{post\_valuation\_status}, \texttt{pct\_acquired}, \texttt{vc\_round},
\texttt{vc\_round\_up\_down\_flat}, \texttt{type\_of\_stock}, \texttt{sellers} \\
Capital structure (deals) &
\texttt{total\_invested\_capital}, \texttt{total\_invested\_equity}, \texttt{debt},
\texttt{total\_new\_debt}, \texttt{total\_debt}, \texttt{raised\_to\_date},
\texttt{total\_raised\_to\_date}, \texttt{total\_investor\_ownership} \\
Financials at deal (deals) &
\texttt{revenue}, \texttt{ebitda}, \texttt{gross\_profit}, \texttt{net\_income},
\texttt{valuation\_ebitda}, \texttt{valuation\_revenue}, \texttt{valuation\_cash\_flow},
\texttt{deal\_size\_ebitda}, \texttt{deal\_size\_revenue},
\texttt{deal\_size\_cash\_flow}, \texttt{revenue\_growth\_since\_last\_deal},
\texttt{number\_of\_employees} \\
Firmography (companies) &
\texttt{description}, \texttt{year\_founded}, \texttt{hq\_location},
\texttt{hq\_country}, \texttt{hq\_region}, \texttt{hq\_sub\_region},
\texttt{employees}, \texttt{financing\_status}, \texttt{business\_status},
\texttt{ownership\_status}, \texttt{company\_universe},
\texttt{total\_money\_raised}, \texttt{last\_financing\_date},
\texttt{last\_financing\_size}, \texttt{last\_financing\_valuation} \\
Investor relation &
\texttt{company\_id}, \texttt{investor\_id}, \texttt{investorname},
\texttt{investorstatus}, \texttt{investorsince} \\
\bottomrule
\end{tabular}
\caption{Field inventory grouped by function. The grouping, not the source table,
determines admissible use: dated deal fields may enter features via the construction of
Section~\ref{sec:data-pit}; snapshot firmography largely may not.}
\label{tab:data-fields}
\end{table}

\section{Deal-type heterogeneity}
\label{sec:data-dealtypes}

The single most consequential schema fact is that \texttt{deal\_class} partitions the
\texttt{deals} table into transaction types that populate \emph{different subsets of
fields}. The main classes are venture capital rounds (angel through late-stage,
identified by \texttt{vc\_round}), M\&A and buyout transactions (control changes,
identified by \texttt{pct\_acquired} and \texttt{sellers}), and debt financings (general
debt, mezzanine, refinancing). A fourth operational bucket --- later-stage growth/PE
minority deals --- behaves as a hybrid of the first two. Treating the table as
homogeneous, e.g.\ pooling \texttt{deal\_size} across classes, mixes a Series~A cheque
with a leveraged buyout price and a term loan, three quantities with different units of
economic meaning. Table~\ref{tab:data-classfields} summarizes typical field population by
class in our synthetic mirror, calibrated to vendor experience.

\begin{table}[htbp]
\centering
\small
\begin{tabular}{p{5.6cm}ccc}
\toprule
Field group & VC round & M\&A / buyout & Debt \\
\midrule
\texttt{deal\_size}, \texttt{deal\_status}          & populated & populated & populated \\
Round structure (\texttt{vc\_round}, up/down/flat,
  \texttt{type\_of\_stock})                          & populated & n.a.      & n.a. \\
\texttt{pre\_money\_valuation}, \texttt{post\_valuation} & sparse & populated & n.a. \\
\texttt{pct\_acquired}, \texttt{sellers}             & sparse    & populated & n.a. \\
Debt fields (\texttt{debt}, \texttt{total\_new\_debt},
  \texttt{total\_debt})                              & n.a.      & populated & populated \\
Financials at deal (\texttt{revenue}, \texttt{ebitda},
  \texttt{net\_income})                              & sparse    & populated & sparse \\
Multiples (\texttt{valuation\_ebitda},
  \texttt{deal\_size\_revenue}, \dots)               & sparse    & populated & n.a. \\
\texttt{total\_investor\_ownership}                  & populated & sparse    & n.a. \\
Cumulative capital (\texttt{raised\_to\_date},
  \texttt{total\_raised\_to\_date})                  & populated & sparse    & sparse \\
\bottomrule
\end{tabular}
\caption{Typical field population by \texttt{deal\_class}. ``Sparse'' means present on a
minority of rows (often mechanically: a seed round has no EBITDA to report); ``n.a.''
means the field is undefined for the class, not missing at random.}
\label{tab:data-classfields}
\end{table}

Three modelling implications follow, and we invoke them repeatedly.

\begin{enumerate}
\item \textbf{Type-specific measurement models.} A mark model for \texttt{deal\_size}
must be fit per class (or with class interactions throughout): VC sizes are
round-indexed and roughly log-normal within round, buyout sizes scale with EBITDA via a
multiple, debt sizes scale with collateral. Pooled moments are meaningless. The same
holds for missingness: an empty \texttt{ebitda} on a seed round is structural (n.a.), on
a buyout it is informative missingness.
\item \textbf{M\&A and IPO are absorbing.} A majority acquisition or listing removes the
firm from the private universe: its counting process $N_i(t)$ stops, not because
intensity fell but because the state space was exited. Deal prediction is therefore a
\emph{competing-risks} problem --- the hazards of ``next VC round'', ``acquired'', and
``fails'' compete for the same firm-time --- previewed here and developed with risk sets
in Chapter~\ref{ch:hawkes}; the selection effects that absorption induces on any
survivors-only analysis are treated in Chapter~\ref{ch:selection}.
\item \textbf{Debt deals inform capital structure, not growth.} A refinancing changes
\texttt{total\_debt} and default risk but carries little signal about revenue growth; a
new venture round carries a valuation mark. In the success models of
Chapters~\ref{ch:success}--\ref{ch:xgb}, debt events enter as capital-structure
covariates, never as success outcomes.
\end{enumerate}

\section{FactSet Standardized Economics}
\label{sec:data-factset}

The macro feed arrives in long format with schema
(\texttt{series\_code}, \texttt{ref\_date}, \texttt{publish\_date}, \texttt{value}):
one row per observation of a series about a reference period, stamped with the date the
figure became public. We use five series (Table~\ref{tab:data-macro-series}). The
distinction between \texttt{ref\_date} and \texttt{publish\_date} is the entire content
of this section: the March CPI describes March but is knowable only in mid-April.

\begin{table}[htbp]
\centering
\small
\begin{tabular}{llll}
\toprule
Code & Series & Frequency & Publication lag \\
\midrule
\texttt{FFUNDT}    & policy target rate       & daily (step function) & same day \\
\texttt{FFUND}     & effective overnight rate & daily                 & next business day \\
\texttt{CPI\_RAW}  & consumer price index, level & monthly            & $\approx 2$ weeks \\
\texttt{CPI\_YOY}  & CPI, year-on-year change & monthly               & $\approx 2$ weeks \\
\texttt{RUNEMP}    & unemployment rate        & monthly               & $\approx 1$ week \\
\bottomrule
\end{tabular}
\caption{Macro series used throughout, FactSet Standardized Economics convention.}
\label{tab:data-macro-series}
\end{table}

\paragraph{As-of join semantics.}
For a series with rows $(r_k, p_k, v_k)$ (reference date, publish date, value), the
\emph{as-of value} at calendar time $t$ is
\begin{equation}
\label{eq:data-asof}
x^{\mathrm{asof}}_t \;=\; v_{k(t)}, \qquad
k(t) \;=\; \arg\max_k \,\{\, r_k : p_k \le t \,\},
\end{equation}
i.e.\ the most recent \emph{reference} among rows already \emph{published}; ties in
$r_k$ (revisions) resolve to the latest $p_k \le t$. Joining on
$r_k \le t$ instead --- the default of most dataframe as-of joins --- silently grants the
model up to two weeks of foresight on CPI, which is material for deal-timing models at
weekly resolution (Chapters~\ref{ch:poisson}--\ref{ch:hawkes}). We also carry the
staleness $t - r_{k(t)}$ as a feature: it is admissible (it is computable at $t$) and
lets the model discount old readings. Figure~\ref{fig:data-macro} shows the synthetic
feed as published.

\begin{figure}[htbp]
\centering
\includegraphics[width=.85\textwidth]{fig_macro}
\caption{The synthetic macro feed \emph{as published}: policy target and effective rate
(daily), CPI year-on-year and unemployment (monthly steps, right-shifted by their
publication lags). The step timing, not the level path, is what a point-in-time model
sees.}
\label{fig:data-macro}
\end{figure}

\section{Point-in-time discipline, formalized}
\label{sec:data-pit}

Let $\mathcal{G}_t$ denote the sigma-algebra generated by all data rows published
strictly before $t$: deal rows with \texttt{deal\_date} $< t$ (subject to the jitter
caveat of Section~\ref{sec:data-quality}), macro rows with \texttt{publish\_date} $< t$,
and static firmography known at first tracking. This is the observational filtration in
the sense of point-process theory \cite{daley2003}, and it is the only information a
deployed model possesses.

\begin{definition}[Admissible feature process]
\label{def:data-admissible}
A feature process $(Z_t)_{t \ge 0}$ is \emph{admissible} if $Z_t$ is
$\mathcal{G}_{t}$-measurable for every $t$, i.e.\ computable from data published before
$t$. A supervised training pair $(Z_{\tau}, Y)$ is admissible if $Z_\tau$ is admissible
at the information date $\tau$ and $Y$ is a function of data published after $\tau$.
\end{definition}

\begin{warning}[Snapshot columns are forbidden]
\label{warn:data-snapshot}
\sloppy
Several \texttt{companies} columns are \emph{end-of-sample snapshots} that reflect the
state of the firm at data delivery, not at any historical $t$:
\texttt{last\_financing\_date}, \texttt{last\_financing\_size},
\texttt{last\_financing\_valuation}, \texttt{total\_money\_raised},
\texttt{financing\_status}, \texttt{business\_status}, and
\texttt{ownership\_status}. Using them as features violates
Definition~\ref{def:data-admissible} at every date before delivery. The leakage is vicious
because it is plausible: \texttt{business\_status} $=$ ``out of business'' predicts
failure spectacularly in backtests and not at all in production. Snapshot columns may be
used only for end-of-sample label construction and for auditing.
\end{warning}

The admissible substitute is \emph{last observation carried forward} (LOCF) over past
deals. For firm $i$ with deal times $t_{i,1} < t_{i,2} < \dots$ and per-deal marks
$m_{i,k}$ (size, valuation, employees, round number),
\begin{equation}
\label{eq:data-locf}
Z_{i,t} \;=\; \Big( m_{i,K_i(t)},\; t - t_{i,K_i(t)},\; K_i(t) \Big),
\qquad K_i(t) = \max\{k : t_{i,k} < t\},
\end{equation}
with an explicit ``no prior deal'' state when $K_i(t) = 0$. The elapsed time
$t - t_{i,K_i(t)}$ is itself a first-class feature: firm state measured at events decays
in relevance, and every model in Parts~II--III uses staleness explicitly. Cumulative
fields (\texttt{raised\_to\_date}) are recomputed from past rows rather than read from
the snapshot, so that the two constructions can be cross-checked for consistency.

\section{Data quality and missingness}
\label{sec:data-quality}

The synthetic mirror injects the defects below at rates calibrated to vendor experience;
all downstream chapters quote robustness against them. Table~\ref{tab:data-quality}
summarizes the defects and which method each one threatens.

\begin{table}[htbp]
\centering
\small
\begin{tabular}{lcp{6.0cm}}
\toprule
Defect & Rate & Main downstream consequence \\
\midrule
\texttt{deal\_size} undisclosed & 28\% & MNAR: small deals hide; mark models need a
missingness indicator (Ch.~\ref{ch:hawkes}) \\
\texttt{deal\_size} estimated only & 15\% & multiplicative error $\approx 18\%$;
errors-in-variables attenuation in size regressions (Ch.~\ref{ch:xgb}) \\
valuation missing & 55\% & valuation-based proxies unusable for half the sample
(Ch.~\ref{ch:us}) \\
sector mislabelled & $\approx 3\%$ & contaminates cross-sector excitation estimates
(Ch.~\ref{ch:hawkes}) \\
duplicate records & $\approx 1.2\%$ & spurious zero-lag self-excitation; inflates
$\kappa$ if not removed \\
date jitter (announced vs.\ closed) & up to 6 days & biases decay-rate estimates;
motivates interval censoring (Ch.~\ref{ch:censored}) \\
out-of-universe deals (no \texttt{companies} row) & $\approx 12\%$ & no firmography
covariates; selection into tracking (Ch.~\ref{ch:selection}) \\
\bottomrule
\end{tabular}
\caption{Defect rates in the synthetic mirror, calibrated to vendor experience, with the
chapter where each consequence is handled.}
\label{tab:data-quality}
\end{table}

Two defects deserve elaboration. First, size missingness is not ignorable: disclosure
correlates negatively with deal size and with the firm being early-stage, so a
complete-case mean overstates typical deal size; every size model carries
$\ind{\texttt{deal\_size\_status} = \text{disclosed}}$ as a covariate and, where the
target is the size itself, a two-part (disclosure, size $\mid$ disclosed) specification.
Second, duplicates and jitter interact badly with self-excitation: a duplicated deal at
lag zero to six days is observationally similar to genuine fast excitation, so
deduplication must precede any Hawkes fit.

\begin{definition}[Deduplication rule]
\label{def:data-dedup}
Two \texttt{deals} rows are merged if they share \texttt{company\_id}, have
\texttt{deal\_date}s within 3 days, and have compatible \texttt{deal\_class}. The
surviving row takes the earlier date and, field by field, the non-missing value (the
larger completeness count breaking conflicts).
\end{definition}

The 3-day window trades off missed duplicates against merged genuine deals: in the
synthetic mirror, it removes 96\% of injected duplicates while merging under 0.1\% of
true distinct pairs, and the residual bias on the branching ratio is below the bootstrap
noise reported in Chapter~\ref{ch:selection}.

\subsection{Entity resolution and identifier hygiene}
\label{sec:data-entity}

Deduplication (Definition~\ref{def:data-dedup}) removes copies of the same transaction;
entity resolution removes copies of the same \emph{firm} or \emph{investor}, and its
failures are worse because they persist across the whole history rather than at one
date. Three rules govern identifier use. First, \texttt{company\_id} is authoritative
but not immortal: vendors occasionally retire an identifier and reissue the firm under
a new one (re-registration, cross-border redomiciliation, relisting after
restructuring). The symptom is an event history that stops dead while a near-namesake
starts from nothing in the same sector and region; the repair is a vendor crosswalk
where one exists, and otherwise a conservative merge requiring simultaneous agreement
on name, country, and \texttt{year\_founded}. An unrepaired split truncates durations
and biases every timing model of Part~III toward short histories. Second, never
\emph{decide} a merge on names. \texttt{company\_name} varies by legal suffix,
stripped diacritics, and renamings; fuzzy matching produces false merges (``Alpha
Capital'' vs.\ ``Alpha Capital Partners'') as readily as false splits. Name similarity
may propose candidate pairs for review; identifiers and structured fields decide.
Third, investors duplicate by design: one economic investor appears as a management
company, as fund vehicles I/II/III, and as regional entities, each with its own
\texttt{investor\_id}. Features that count distinct investors or measure investor
experience roll vehicles up to the management company; any analysis of co-investment
structure must state its rollup, because the rollup changes the network. Finally,
referential integrity is checked mechanically at every build: each
\texttt{deals.company\_id} resolves in \texttt{companies} or is flagged
out-of-universe (the $\approx 12\%$ of Table~\ref{tab:data-quality}); each
\texttt{investor\_id} in the investor relation resolves; and \texttt{investorsince}
never predates \texttt{year\_founded}. Violating rows are quarantined rather than
silently dropped, so that their rate can be monitored across vendor deliveries.

\section{Worked example: constructing the weekly modelling panel}
\label{sec:data-worked}

The constructions of this chapter are best fixed by running them once, end to end, on
a small window. The target is the weekly panel consumed by
Chapters~\ref{ch:poisson}--\ref{ch:hawkes}: on a weekly grid $t_1 < t_2 < \dots$
(Friday closes), a firm-week feature array, an activity mask, and an event list. Five
passes produce these from the raw tables; each pass exists to prevent a specific
failure, named as we go.

\begin{enumerate}
\item \textbf{Deduplication pass.} Apply Definition~\ref{def:data-dedup} to
\texttt{deals} before anything else reads the table, logging row counts before and
after. \emph{Failure prevented:} duplicate rows at lags of 0--6 days masquerade as
fast self-excitation and inflate the branching ratio
(Section~\ref{sec:data-quality}); they also double-count capital in the cumulative
fields recomputed in pass 3.
\item \textbf{As-of macro join.} For each grid date $t$ and each series of
Table~\ref{tab:data-macro-series}, attach the as-of value (\ref{eq:data-asof}) and
the staleness $t - r_{k(t)}$. Table~\ref{tab:data-asof-example} traces
\texttt{CPI\_YOY} through three consecutive grid dates around the April publication:
the March figure (reference 2019-03-31, published 2019-04-10) is invisible on
2019-04-05 and governs the two later weeks with growing staleness.
\emph{Failure prevented:} joining on \texttt{ref\_date} attaches the March print to
the 2019-04-05 row --- five days of foresight on this series, up to two weeks in
general.
\item \textbf{LOCF firm panel.} For each firm $i$ and grid date $t$, build $Z_{i,t}$
of (\ref{eq:data-locf}): the marks of the last deal strictly before $t$, the elapsed
time since it, the deal count $K_i(t)$, and cumulative capital recomputed by summing
past deal rows, with an explicit no-prior-deal state when $K_i(t) = 0$.
\emph{Failure prevented:} reading firm state from the \texttt{companies} snapshot
columns leaks end-of-sample status (Warning~\ref{warn:data-snapshot}); LOCF over
dated rows is the admissible substitute.
\item \textbf{Active mask.} Set $a_{i,t} = 1$ iff firm $i$ is tracked and unabsorbed
at $t$: its first observed deal (or vendor add date) precedes $t$ and no absorbing
M\&A, IPO, or failure event occurred before $t$
(Section~\ref{sec:data-dealtypes}). \emph{Failure prevented:} pre-tracking weeks
pretend the firm was at risk while it was invisible, and post-absorption weeks count
firm-time in which no event could occur; both distort intensity estimates, which is
the risk-set discipline formalized in Chapters~\ref{ch:hawkes}
and~\ref{ch:selection}.
\item \textbf{Event list.} Emit the surviving deals as
(\texttt{company\_id}, \texttt{deal\_date}, \texttt{deal\_class}, marks), restricted
to the classes the model in question treats as events, with debt deals routed to
covariates rather than outcomes. \emph{Failure prevented:} pooling refinancings with
venture rounds mixes units of economic meaning
(Section~\ref{sec:data-dealtypes}).
\end{enumerate}

\begin{table}[htbp]
\centering
\small
\begin{tabular}{lllcc}
\toprule
Grid date $t$ & Latest publication $\le t$ & Reference month & Value (\%) & Staleness (days) \\
\midrule
2019-04-05 & 2019-03-13 & 2019-02 & 1.9 & 36 \\
2019-04-12 & 2019-04-10 & 2019-03 & 2.1 & 12 \\
2019-04-19 & 2019-04-10 & 2019-03 & 2.1 & 19 \\
\bottomrule
\end{tabular}
\caption{Three weeks of the as-of join (\ref{eq:data-asof}) for \texttt{CPI\_YOY}.
The publication lag makes the February figure the newest \emph{knowable} reading on
2019-04-05 although March has ended; the staleness $t - r_{k(t)}$ resets at each
publication and then grows by 7 days per week.}
\label{tab:data-asof-example}
\end{table}

The output is three aligned objects --- the feature array over firm-weeks with
$a_{i,t} = 1$, the mask itself, and the event list --- and nothing else crosses into
the modelling code. Before release, the following audit runs mechanically; any failed
item blocks the build.

\begin{enumerate}
\item \textbf{Dedup accounting.} The merge rate of pass 1 is compared with the
expected $\approx 1.2\%$ of Table~\ref{tab:data-quality}; a large deviation signals a
changed vendor delivery, not a modelling fact.
\item \textbf{No macro foresight.} Every joined macro row satisfies
\texttt{publish\_date} $< t$; the assertion runs on the joined output, not on trust
in the join code.
\item \textbf{LOCF causality.} For a random sample of firm-weeks, delete all raw rows
dated on or after $t$ and rebuild $Z_{i,t}$; the rebuilt values must be identical.
\item \textbf{Mask consistency.} Every event falls in a firm-week with
$a_{i,t} = 1$, and no firm has $a_{i,t} = 1$ after an absorbing event.
\item \textbf{Snapshot blocklist.} No column of Warning~\ref{warn:data-snapshot}
appears in the feature array, enforced by name against a schema blocklist; the
recomputed cumulative capital is reconciled against the end-of-sample snapshot, with
discrepancies logged.
\end{enumerate}

\section{EU versus US coverage: panel versus events}
\label{sec:data-euus}

The schemas above are common to both regimes; what differs is what can be attached to
them. For European firms, statutory registry filings supply a yearly fundamentals panel
$F_{i,y}$ with near-complete coverage --- an \emph{observation of the state between
deals}. For US firms, the deal rows of this chapter are essentially all there is: state
is measured only at the firm's own transactions, through the sparse financials-at-deal
group of Table~\ref{tab:data-fields}. This is the axis of the manual announced in
Chapter~\ref{ch:problem}. Part~I (Chapters~\ref{ch:success}--\ref{ch:causal}) exploits
the EU panel for supervised and causal analysis of success; Part~II
(Chapter~\ref{ch:us}) asks what the deal marks alone can proxy; Part~III
(Chapters~\ref{ch:poisson}--\ref{ch:selection}) embraces the event-only view and models
the deal stream itself. All three parts share the point-in-time constructions
(\ref{eq:data-asof}) and (\ref{eq:data-locf}) verbatim.
